\section{Parameter Shift Rule (PSR)}\label{appendix-f}

\emph{Parameter Shift Rule} adalah teknik fundamental dalam komputasi kuantum variabel (Variational Quantum Circuits) untuk menghitung gradien eksak dari nilai ekspektasi terhadap parameter sirkuit. Berbeda dengan \emph{finite difference} yang hanya memberikan estimasi, PSR memberikan nilai gradien yang analitik dan eksak.

\hypertarget{sejarah-singkat}{%
\subsection{Sejarah Singkat}\label{sejarah-singkat}}

Konsep \emph{Parameter Shift Rule} pertama kali diperkenalkan ke komunitas \emph{Quantum Machine Learning} oleh \textbf{Mitarai et al.} dalam makalah berjudul \emph{``Quantum Circuit Learning''} pada tahun 2018. Tak lama kemudian, \textbf{Schuld et al.} memperluas dan menggeneralisasi metode ini dalam makalah \emph{``Evaluating analytic gradients of quantum circuits on quantum hardware''}.

Sebelum adanya PSR, peneliti sering menggunakan metode \emph{finite difference} (seperti metode Euler) yang memiliki kelemahan pada \emph{hardware} kuantum yang berisik karena ketidakstabilan numerik saat \(\Delta \theta\) sangat kecil. PSR memecahkan masalah ini dengan menggunakan pergeseran parameter yang besar (misal \(\pi/2\)).

\hypertarget{definisi-matematis}{%
\subsection{Definisi Matematis}\label{definisi-matematis}}

Misalkan kita memiliki fungsi energi (nilai ekspektasi) \(E(\theta)\): \[E(\theta) = \langle \psi | U^\dagger(\theta) M U(\theta) | \psi \rangle\]

Di mana \(U(\theta)\) adalah gerbang unitari yang dibangkitkan oleh operator Pauli \(\sigma\) (seperti \(X, Y, Z\)): \[U(\theta) = \exp\left(-i \frac{\theta}{2} \sigma\right)\]

Persamaan \emph{Parameter Shift Rule} menyatakan bahwa turunan \(E(\theta)\) terhadap \(\theta\) adalah: \[\frac{\partial E(\theta)}{\partial \theta} = \frac{1}{2} \left[ E\left(\theta + \frac{\pi}{2}\right) - E\left(\theta - \frac{\pi}{2}\right) \right]\]

\hypertarget{penurunan-rumus-derivation}{%
\subsection{Penurunan Rumus (Derivation)}\label{penurunan-rumus-derivation}}

Mari kita turunkan secara bertahap menggunakan sifat-sifat matriks Pauli.

\hypertarget{a.-turunan-langsung}{%
\subsubsection{Turunan Langsung}\label{a.-turunan-langsung}}

Gunakan aturan perkalian (product rule) pada \(E(\theta)\): \[\frac{\partial E}{\partial \theta} = \langle \psi | \frac{\partial U^\dagger}{\partial \theta} M U | \psi \rangle + \langle \psi | U^\dagger M \frac{\partial U}{\partial \theta} | \psi \rangle\]

Karena \(U(\theta) = e^{-i\frac{\theta}{2}\sigma}\), maka: \[\frac{\partial U}{\partial \theta} = -\frac{i}{2}\sigma U(\theta), \quad \frac{\partial U^\dagger}{\partial \theta} = \frac{i}{2} U^\dagger(\theta) \sigma\]

Substitusikan ke persamaan turunan: \[\frac{\partial E}{\partial \theta} = \langle \psi | \left(\frac{i}{2} U^\dagger \sigma\right) M U | \psi \rangle + \langle \psi | U^\dagger M \left(-\frac{i}{2} \sigma U\right) | \psi \rangle\] \[\frac{\partial E}{\partial \theta} = \frac{i}{2} \langle \psi(\theta) | \sigma M - M \sigma | \psi(\theta) \rangle\] \[\frac{\partial E}{\partial \theta} = \frac{i}{2} \langle \psi(\theta) | [\sigma, M] | \psi(\theta) \rangle \quad \dots (1)\]

\hypertarget{b.-ekspansi-pergeseran-parameter}{%
\subsubsection{Ekspansi Pergeseran Parameter}\label{b.-ekspansi-pergeseran-parameter}}

Gunakan identitas Euler untuk operator Pauli (\(\sigma^2 = I\)): \[U(s) = e^{-i\frac{s}{2}\sigma} = \cos\left(\frac{s}{2}\right) I - i \sin\left(\frac{s}{2}\right) \sigma\]

Evaluasi \(E(\theta + s)\) di mana \(s\) adalah pergeseran: \[E(\theta + s) = \langle \psi(\theta) | U^\dagger(s) M U(s) | \psi(\theta) \rangle\] \texttt{mengapa\ s\ hanya\ berpengaruh\ pada\ U\ saja\ dan\ tidak\ pada\ \$\textbackslash{}psi\$} \textgreater{} \textbf{Penjelasan:} Hal ini didasarkan pada sifat eksponensial operator unitari. Karena \(U(\theta + s) = e^{-i\frac{\theta+s}{2}\sigma} = e^{-i\frac{s}{2}\sigma} e^{-i\frac{\theta}{2}\sigma} = U(s)U(\theta)\), maka kita dapat menuliskan: \textgreater{} \[E(\theta + s) = \langle \psi | U^\dagger(\theta) U^\dagger(s) M U(s) U(\theta) | \psi \rangle\] \textgreater{} Dengan mendefinisikan \(|\psi(\theta)\rangle = U(\theta)|\psi\rangle\), persamaan di atas menjadi: \textgreater{} \[E(\theta + s) = \langle \psi(\theta) | U^\dagger(s) M U(s) | \psi(\theta) \rangle\] \textgreater{} Jadi, \(s\) dianggap sebagai pergeseran lokal (rotasi tambahan) pada state yang sudah dipreparasi oleh parameter \(\theta\).

\[E(\theta + s) = \langle \psi(\theta) | (\cos\frac{s}{2} I + i \sin\frac{s}{2} \sigma) M (\cos\frac{s}{2} I - i \sin\frac{s}{2} \sigma) | \psi(\theta) \rangle\]

Setelah melakukan ekspansi aljabar: \[E(\theta + s) = \cos^2\frac{s}{2} \langle M \rangle + \sin^2\frac{s}{2} \langle \sigma M \sigma \rangle + \cos\frac{s}{2}\sin\frac{s}{2} \cdot i \langle [\sigma, M] \rangle\]

Dengan cara yang sama untuk \(E(\theta - s)\): \[E(\theta - s) = \cos^2\frac{s}{2} \langle M \rangle + \sin^2\frac{s}{2} \langle \sigma M \sigma \rangle - \cos\frac{s}{2}\sin\frac{s}{2} \cdot i \langle [\sigma, M] \rangle\]

\hypertarget{c.-finalisasi}{%
\subsubsection{Finalisasi}\label{c.-finalisasi}}

Kurangi kedua persamaan tersebut: \[E(\theta + s) - E(\theta - s) = 2 \cos\frac{s}{2} \sin\frac{s}{2} \cdot i \langle [\sigma, M] \rangle\] Gunakan identitas trigonometri \(2\sin A \cos A = \sin 2A\): \[E(\theta + s) - E(\theta - s) = \sin(s) \cdot i \langle [\sigma, M] \rangle \quad \dots (2)\]

Bandingkan persamaan (2) with persamaan (1). Jika kita pilih \(s = \pi/2\): \[\sin(\pi/2) = 1\] Maka: \[E(\theta + \pi/2) - E(\theta - \pi/2) = i \langle [\sigma, M] \rangle\]

Karena dari persamaan (1) kita tahu \(\frac{\partial E}{\partial \theta} = \frac{1}{2} \left( i \langle [\sigma, M] \rangle \right)\), maka terbukti bahwa: \[\frac{\partial E}{\partial \theta} = \frac{1}{2} \left[ E\left(\theta + \frac{\pi}{2}\right) - E\left(\theta - \frac{\pi}{2}\right) \right]\]

\hypertarget{keunggulan-psr}{%
\subsection{Keunggulan PSR}\label{keunggulan-psr}}

\begin{enumerate}
\def\labelenumi{\arabic{enumi}.}
\tightlist
\item
  \textbf{Analitik}: Memberikan gradien yang benar-benar eksak, bukan pendekatan.
\item
  \textbf{Noise Robust}: Karena pergeseran \(s\) cukup besar (\(\pi/2\)), sinyal yang dihasilkan lebih kuat dibandingkan \emph{noise} \emph{hardware}, berbeda dengan \emph{finite difference} yang butuh \(s\) sangat kecil.
\item
  \textbf{Implementasi Mudah}: Tidak memerlukan sirkuit tambahan atau tambahan \emph{qubit}; cukup menjalankan sirkuit yang sama dengan parameter yang berbeda.
\end{enumerate}
